% Materiał z AF 25; fragment do rozdziału o dystrybucjach.
\subsection{Lokalność i problemy przedłużania dystrybucji}

\begin{twierdzenie}[Ograniczenie i brak automatycznego przedłużenia]
\label{thm:restriction25}
Jeżeli $U\subseteq\Omega\subseteq\R^d$ są niepustymi
zbiorami otwartymi, to odwzorowanie ograniczenia
$\mathcal D'(\Omega)\to\mathcal D'(U)$ jest liniowe
i ciągłe w topologiach słabych-* oraz silnych.
Nie dla każdej pary $U\subseteq\Omega$ odwzorowanie
ograniczenia jest surjektywne. Nawet dla przedziałów
$U=(0,1)$ i $\Omega=(-1,1)$ istnieją dystrybucje
regularne na $U$, które nie mają przedłużenia na $\Omega$.
\end{twierdzenie}
\begin{proof}
Funkcję testową $\varphi\in\mathcal D(U)$ przedłużamy
przez zero do $\Omega$ i oznaczamy przedłużenie przez
$j\varphi=\widetilde\varphi$. Jeżeli
$\supp\varphi\subseteq L\Subset U$, to $L$ jest również
zwarty w $\Omega$, a każdy punkt $\Omega\setminus U$
ma otoczenie rozłączne z $L$. Przedłużenie jest więc
gładkie także na brzegu $U$ w $\Omega$ i
$\supp(j\varphi)=\supp\varphi$, gdzie oba nośniki
utożsamiamy jako podzbiory $\Omega$.
Na każdym niezerowym etapie $\mathcal D_L(U)$
odwzorowanie $j$ ma wartości w $\mathcal D_L(\Omega)$
i zachowuje seminormy $p_{L,m}$, ponieważ pochodne
na $U$ nie zmieniają się. Na etapie zerowym ciągłość
jest oczywista. Własność uniwersalna granicy induktywnej
(Twierdzenie~\ref{thm:LF-universal18}) zapewnia ciągłość
liniowego odwzorowania $j:\mathcal D(U)\to\mathcal D(\Omega)$.

Ograniczenie $R_U T=T|_U$ definiujemy przez
$\langle R_U T,\varphi\rangle=\langle T,j\varphi\rangle$.
Jest to ciągły funkcjonał na $\mathcal D(U)$ jako
złożenie ciągłych odwzorowań liniowych.
Dla każdego ograniczonego $B\subseteq\mathcal D(U)$
mamy dokładną równość
\[
 \sup_{\varphi\in B}|\langle R_U T,\varphi\rangle|
 =\sup_{\eta\in j(B)}|\langle T,\eta\rangle|.
\]
Ponieważ $j(B)$ jest ograniczony, otrzymujemy ciągłość
silną, a dla $B$ jednoelementowego także słabą-*.
Nie zakładamy tu gęstości obrazu $j$.

Dla otwartego $V\subseteq U$ kolejne przedłużenia przez
zero dają to samo odwzorowanie co bezpośrednie
przedłużenie z $V$ do $\Omega$. Zatem
$(T|_U)|_V=T|_V$. Z definicji nośnika i
Lematu~\ref{lem:distribution-local19} wynika
$\supp(T|_U)=U\cap\supp T$, z nośnikiem po lewej
stronie rozumianym względnie w $U$: punkt $x\in U$
nie należy do tego nośnika wtedy i tylko wtedy,
gdy istnieje otwarte otoczenie $V\subseteq U$ punktu
$x$, na którym $T|_V=0$. W szczególności
$\ker R_U=\{T\in\mathcal D'(\Omega):
\supp T\subseteq\Omega\setminus U\}$.
Jeżeli przedłużenie istnieje, jego jednoznaczność
nie jest na ogół zapewniona: do dowolnego przedłużenia
można dodać element tego jądra.

Aby wykazać drugie stwierdzenie, weźmy
$U=(0,1)$, $\Omega=(-1,1)$ i dystrybucję regularną
$T_f$ wyznaczoną przez $f(x)=\exp(1/x^2)$ na $U$.
Funkcja $f$ jest gładka na $U$, a więc lokalnie
całkowalna. Ustalmy $K=[-1/2,1/2]\Subset\Omega$.
Gdyby istniało przedłużenie
$\widetilde T\in\mathcal D'(\Omega)$ z
$\widetilde T|_U=T_f$, to z
Twierdzenia~\ref{thm:distribution-estimate19}
dla tego ustalonego $K$ istniałyby $m\in\N_0$
i $C>0$ takie, że
$|\langle\widetilde T,\varphi\rangle|\le C p_{K,m}(\varphi)$
dla wszystkich $\varphi\in\mathcal D_K(\Omega)$.
Niech $\psi\in\mathcal D((-1,1))$ będzie nieujemna
i niezerowa. Przedłużamy ją przez zero na $\R$.
Dla całkowitych $n\ge3$ przyjmujemy
$r_n=1/(10n^2)$ i
$\varphi_n(x)=\psi((x-1/n)/r_n)$.
Nośnik zawiera się w przedziale
$[1/n-r_n,1/n+r_n]\subset(0,1/2)$, zatem
$\varphi_n$ jest testem zarówno na $U$, jak i na
ustalonym etapie $\mathcal D_K(\Omega)$.
Dla każdego $0\le k\le m$ mamy
$\varphi_n^{(k)}(x)=r_n^{-k}\psi^{(k)}((x-1/n)/r_n)$.
Ponieważ $r_n\le1$, otrzymujemy
\[
 p_{K,m}(\varphi_n)
 \le 10^m n^{2m}\max_{0\le k\le m}
 \sup_{\R}|\psi^{(k)}|.
\]
Z drugiej strony dla $x=1/n+r_nt$, $|t|\le1$,
mamy $0<x\le(11/10)n^{-1}$, a więc
$1/x^2\ge(100/121)n^2\ge n^2/2$.
Podstawienie w całce i nieujemność $\psi$ dają
\[
 \begin{aligned}
 \langle T_f,\varphi_n\rangle
 &=r_n\int_{-1}^1 e^{1/(1/n+r_nt)^2}\psi(t)\,
 \mathrm d\lambda_1(t)\\
 &\ge \frac{e^{n^2/2}}{10n^2}
 \int_{-1}^1\psi(t)\,\mathrm d\lambda_1(t).
 \end{aligned}
\]
Stała $c=\frac1{10}\int_{-1}^1\psi(t)\,\mathrm d\lambda_1(t)$
jest ściśle dodatnia. Ponieważ
$\langle\widetilde T,\varphi_n\rangle=
\langle T_f,\varphi_n\rangle$, z oszacowania na
$\mathcal D_K(\Omega)$ wynikałoby
$c e^{n^2/2}\le C' n^{2m+2}$ dla wszystkich $n\ge3$,
gdzie $C'$ nie zależy od $n$. Jednak
$e^{n^2/2}\ge (n^2/2)^{m+2}/(m+2)!$,
więc po podzieleniu przez $n^{2m+2}$ lewa strona
rośnie co najmniej jak dodatnia stała razy $n^2$.
Otrzymujemy sprzeczność. Brak przedłużenia dotyczy
wszystkich dystrybucji, nie tylko dystrybucji regularnych.
\end{proof}
